\documentclass[12pt]{article}



\usepackage{amsmath}
\usepackage{amssymb}
\usepackage{amsthm}


\begin{document}
\begin{enumerate}
	\item 
		\begin{itemize}
			\item[(a)] We want to show that $[0,1] \subseteq \mathbb{R}$ endowed with the subspace topology is connected. Assume it's disconnected, so there exists $U,V \subseteq [0,1]$ open such that $U \cap V = \emptyset$ and $U \cup V = [0,1]$. Note that all the open sets in $[0,1]$ are generated from the set $G = \{[0,a), (b,1], (c,d) \colon a \in (0,1], b \in [0,1), c,d \in [0,1]\}$. Suppose $1 \in V$ WLOG, and let $s = \sup U \cap [0,1]$, $s$ cannot be in $U$ since it cannot be one, and since $U$ is open, there exists an open ball of radius $\varepsilon$ that contains $u$, so there exists a greater element than $s$ in the open ball that is in $U$. $s$ cannot belong to $V$ either, it cannot be $1$ since that would imply that $V$ is not open (it would contain one single point), and it cannot be $0$ either since $U$ wouldn't be open, so $0 < s < 1$, but since $V$ open $s$ is contained in an open ball $B \subseteq V$, so we can find an element of $V$ smaller than $s$ that would still be an upper bound for $U \cap [0,1]$. 
			\item[(b)] Suppose $\bar{A}$ is not connected, that means there exists $U,V$ open in $\bar{A}$ that partition $\bar{A}$.  But then 
				$$A = (U \cap A) \cup (V \cap A)$$. $U \cap A$
				and $V \cap A$ are open in $A \subseteq \bar{A}$, that means one of the two is empty. Suppose $U \cap A = \emptyset$, then $$V \cap A = A \implies V =A$$. That means $U = \bar{A} \setminus A$,which is not an open set, since letting $x \in \bar{A} \setminus A = \bar{A} \cap A^c$, any open neighborhood of $x$ will intersect $A^c$.
			\item Suppose $f(X)$ not connected, then there exists $U,V$ open in $f(X)$ such that $U \cap V = \emptyset$, and $U \cup V = f(X)$, then applying the preimage, we have 
				$$
				f^{-1}(U \cap V) = f^{-1}(U) \cap f^{-1}(V) = \emptyset
				$$
				and 
				$$
				f^{-1}(U) \cup f^{-1}(V) = X
				$$
				Since $f^{-1}(U)$ and $f^{-1}(V)$ are open in $X$, this leads to a contradiction since those would partition $X$, a connected set. 
			
		\end{itemize}
	\item $(\implies)$ let $F \subseteq Y$ be a closed set in $Y$, then $F^c$ is open, that means 
		$$
			f^{-1}(F^c) = f^{-1}(F)^c 
		$$
		is open in $X$ so $f^{-1}(F)$ is closed in $X$\\
		Conversely, let $U$ be an open set in $Y$, then $f^{-1}(U^c) = f^{-1}(U)^c$ is closed in $X$, that means $f^{-1}(U)$ is open in $X$. 
	\item \begin{itemize}
			\item[(a)] Letting $W$ be an open set in $Y$
		\begin{align*}
			f^{-1}(W) &= (f^{-1}(W) \cap A) \cup (f^{-1}(W) \cap B)\\
				  &= \overbrace{f_{| A}^{-1}(W)}^{\text{open in $X$}} \cup \overbrace{f_{| B}^{-1}(W)}^{\text{open in $X$}} \\
		\end{align*}
	\item[(b)]
		Using question number 2, letting $F$ be a closed set in $Y$
		\begin{align*}
			f^{-1}(F) &= (f^{-1}(F) \cap A) \cup (f^{-1}(F) \cap B)\\
				  &= \overbrace{f_{| A}^{-1}(F)}^{\text{closed  in $X$}} \cup \overbrace{f_{| B}^{-1}(F)}^{\text{closed in $X$}} \\
		\end{align*}
	\item IDK

		\end{itemize}
	\item We know that $$\Delta^c = \{(x,y) \colon x\neq y, x \in X, y \in Y\}$$ is open in $X \times X$. Thus letting $x,y \in X$ and $x \neq y$, we know that $(x,y) \in \Delta^c$, so there exists a open neighborhood $U \times V \subseteq \Delta^c$ that is open in the product topology of $X \times X$ such that $x \in U$ and $y \in V$, and by definition of $\Delta^c$, $U \cap V = \emptyset$, which gives us the necessary condition for Hausdorff.
\end{enumerate}

\end{document}
